\chapter{Experiments with Heston model dynamics}
\label{chap:heston}

This chapter provides experimental results  of the SRForgan applied to Heston model trajectories. 
In particular the \texttt{HestonSimulator}, a data-generation object that produces trajectories of financial log-prices under the Heston stochastic volatility model and computes the conditional probability density function of future log-prices using the Fractional Fast Fourier Transform (frFFT).
Every mathematical concept required to understand the implementation is introduced from first principles.

% ─────────────────────────────────────────────────────────────────────────────
\section{The Heston Stochastic Volatility Model}
\label{sec:heston_model}

The classical Black--Scholes model assumes that the instantaneous volatility of an asset price is a deterministic constant $\sigma$.
This assumption leads to a log-normal distribution for the terminal asset price, which is analytically tractable but empirically inadequate since financial markets exhibit a so-called
\textit{volatility smile}, whereby options at different strikes imply different values of $\sigma$ 
when inverted through the Black--Scholes formula. 
This observation alone establishes that volatility cannot be constant.

\subsection{The Heston System of SDEs}

To remedy this shortcoming, Heston \cite{heston:heston_paper} introduced in 1993 a two-factor model in which both the log-price $X_t = \ln S_t$ and the instantaneous variance $v_t$ are stochastic.
The system of stochastic differential equations (SDEs) in the \texttt{HestonSimulator} is written in \textit{log-price form}:

\begin{align}
    &dX_t = \left(\mu - \tfrac{1}{2}v_t\right)dt + \sqrt{v_t}\, dW_t^X \label{eq:logprice_sde_heston}\\
    &dv_t = \kappa(\theta - v_t)\,dt + \sigma_v\sqrt{v_t}\, dW_t^v \label{eq:var_sde_heston}\\
    &\langle dW_t^X,\, dW_t^v \rangle = \rho\, dt \label{eq:correlation_heston}
\end{align}

Each term carries a precise financial interpretation.
\begin{itemize}
    \item Log-price SDE (Eq. \ref{eq:logprice_sde_heston}). Applying It\^{o}'s lemma to $X_t = \ln S_t$ under the assumption that $S_t$ follows a geometric Brownian motion with stochastic diffusion coefficient $\sqrt{v_t}$ yields the drift $\mu - v_t/2$. The term $\mu$ is the risk-neutral drift (typically the risk-free rate), while the $-v_t/2$ correction, 
    called the \textit{It\^{o} correction}, arises from the quadratic variation of $X_t$ and ensures that $\mathbb{E}[S_t] = S_0 e^{\mu t}$.
    \item Variance SDE (Eq.\,\ref{eq:var_sde_heston}). The variance $v_t$ evolves according to a \textit{Cox--Ingersoll--Ross} (CIR) process \cite{heston:cir_process}.
     The three parameters have well-defined roles: $\kappa > 0$ is the \textit{mean-reversion speed}, 
     controlling how quickly $v_t$ is pulled back toward its long-run equilibrium; $\theta > 0$ is
      the \textit{long-run variance}, the level to which $v_t$ reverts; and $\sigma_v > 0$ is the \textit{volatility of volatility} (vol-of-vol), 
      governing the magnitude of random fluctuations in $v_t$.
      \item Correlation (Eq.\,\ref{eq:correlation_heston}). The two Brownian motions $W^X$ and $W^v$ are correlated with instantaneous correlation $\rho \in (-1, 1)$. For equities, $\rho$ is typically negative: a negative shock to the asset price tends to be accompanied by an increase in variance, a well-documented empirical phenomenon known as the \textit{leverage effect}.
\end{itemize}

A key requirement for the CIR process is that $v_t \geq 0$ at all times. it has been proved that the boundary $v_t = 0$ is never reached (i.e., $v_t > 0$ almost surely) if and only if

\begin{equation}
    2\kappa\theta > \sigma_v^2.
    \label{eq:feller}
\end{equation}

When condition the Feller condition (\ref{eq:feller}) is violated, the process touches zero with positive probability. In the simulator, this boundary is handled numerically by the \textit{full-truncation} scheme described in Section~\ref{sec:discretisation}.\\

% ─────────────────────────────────────────────────────────────────────────────

To simulate and compute the correlated pair $(dW_t^X, dW_t^v)$ over a discrete time step $\Delta t$, one exploits the fact that any bivariate Gaussian vector with correlation $\rho$ can be constructed from two independent standard normals $Z_1, Z_2 \sim \mathcal{N}(0,1)$ via the \textit{Cholesky decomposition}. Specifically, the increments

\begin{align}
    \Delta W^X &= \sqrt{\Delta t}\, Z_1 \label{eq:dWX}\\
    \Delta W^v &= \sqrt{\Delta t}\left(\rho\, Z_1 + \sqrt{1-\rho^2}\, Z_2\right) \label{eq:dWv}
\end{align}

satisfy $\mathbb{E}[\Delta W^X] = \mathbb{E}[\Delta W^v] = 0$, $\text{Var}(\Delta W^X) = \text{Var}(\Delta W^v) = \Delta t$, and $\text{Cov}(\Delta W^X, \Delta W^v) = \rho\,\Delta t$, which matches the covariance structure prescribed by Eq.~(\ref{eq:correlation_heston}). 
% ─────────────────────────────────────────────────────────────────────────────
\subsection{Numerical Discretisation of the Heston SDEs}
\label{sec:discretisation}

Since equations (\ref{eq:logprice_sde_heston})--(\ref{eq:var_sde_heston}) do not admit closed-form solutions in general, they must be integrated numerically. The proposed simulator supports two schemes:\\
The \textit{Euler--Maruyama} scheme \cite{heston:euler_maruyama} is the SDE analogue of Euler's method for ODEs. For a generic It\^{o} SDE $dY_t = a(Y_t)\,dt + b(Y_t)\,dW_t$, the scheme reads

\begin{equation}
    Y_{n+1} = Y_n + a(Y_n)\,\Delta t + b(Y_n)\,\Delta W_n.
    \label{eq:euler}
\end{equation}

Applied to the log-price SDE (\ref{eq:logprice_sde_heston}), with $\hat{v}_n = \max(v_n, 0)$ (the truncated variance):

\begin{equation}
    X_{n+1} = X_n + \left(\mu - \tfrac{1}{2}\hat{v}_n\right)\Delta t + \sqrt{\hat{v}_n}\,\Delta W_n^X.
    \label{eq:euler_logprice}
\end{equation}

The Euler--Maruyama scheme has strong convergence order $1/2$ and weak convergence order $1$ for SDEs with globally Lipschitz coefficients.
The second proposed discretization technique is the \textit{Milstein scheme} \cite{heston:milstein} improves upon Euler--Maruyama by including an additional correction term derived from It\^{o}'s formula applied to the diffusion coefficient. For an SDE with diffusion $b(Y) = \sigma_v\sqrt{Y}$, the correction is

\begin{equation}
    \frac{1}{2}b(Y_n)\frac{\partial b}{\partial Y}(Y_n)\bigl((\Delta W_n)^2 - \Delta t\bigr) = \frac{\sigma_v^2}{4}\bigl((\Delta W_n^v)^2 - \Delta t\bigr).
    \label{eq:milstein_correction}
\end{equation}

This term captures the curvature of the diffusion coefficient and raises the strong convergence order from $1/2$ to $1$.\\

Because the CIR process can reach zero when the Feller condition fails (or nearly so), a naive Milstein step may produce $v_{n+1} < 0$, which would make $\sqrt{v_{n+1}}$ undefined at the next step. A \textit{full-truncation} strategy addresses this in two stages:

\begin{enumerate}
    \item Before computing the drift and diffusion, replace $v_n$ with $\hat{v}_n = \max(v_n, 0)$ inside all coefficient functions.
    \item After the update step, apply a second truncation: $v_{n+1} \leftarrow \max(v_{n+1}, 0)$.
\end{enumerate}

The complete full-truncation Milstein step for the variance is therefore

\begin{equation}
    v_{n+1} = \max\!\Bigl(v_n + \kappa(\theta - \hat{v}_n)\Delta t + \sigma_v\sqrt{\hat{v}_n}\,\Delta W_n^v + \tfrac{1}{4}\sigma_v^2\bigl((\Delta W_n^v)^2 - \Delta t\bigr),\; 0\Bigr).
    \label{eq:milstein_full_trajectory}
\end{equation}


% ─────────────────────────────────────────────────────────────────────────────
\subsection{The Characteristic Function and the Fourier Inversion Formula}

The \textit{characteristic function} of a real-valued random variable $X$ is defined as

\begin{equation}
    \varphi(u) = \mathbb{E}\!\left[e^{iuX}\right], \quad u \in \mathbb{R},
    \label{eq:cf_def}
\end{equation}

where $i = \sqrt{-1}$. The Characteristic Function is the Fourier transform of the PDF and encodes all information about the distribution
 of $X$. Its fundamental importance stems from the fact that many stochastic models, including the Heston model, 
 admit a closed-form CF even when the PDF itself does not exist in closed form. If $\varphi$ is known analytically, the PDF can be recovered numerically via Fourier inversion.\\
For the Heston model, \cite{heston:heston_paper} it has been derived the conditional CF of the log-price $X_{t+\tau}$ given the state $(X_t, v_t)$ at time $t$.
The CF takes the affine-exponential form

\begin{equation}
    \varphi(u;\tau) = \exp\!\Bigl(iuX_0 + iu\mu\tau + A(u,\tau) + B(u,\tau)\,v_0\Bigr),
    \label{eq:cf_heston}
\end{equation}

where $X_0 = X_t$ is the current log-price, $v_0 = v_t$ is the current variance, and $\tau$ is the forecast horizon. The function $A$ captures the long-run contribution while $B$ captures the effect of the current variance. They both depend on the model parameters through an intermediate quantity $d$:

\begin{align}
    d &= \sqrt{(\kappa - \rho\sigma_v\, iu)^2 + \sigma_v^2(u^2 + iu)}, \quad \text{Re}(d) \geq 0, \label{eq:d_heston}\\
    \tilde{g} &= \frac{\kappa - \rho\sigma_v\, iu - d}{\kappa - \rho\sigma_v\, iu + d}, \label{eq:g_tilde}\\
    A(u,\tau) &= \frac{\kappa\theta}{\sigma_v^2}\left[(\kappa - \rho\sigma_v\, iu - d)\tau - 2\ln\!\frac{1 - \tilde{g}\,e^{-d\tau}}{1 - \tilde{g}}\right], \label{eq:A_heston}\\
    B(u,\tau) &= \frac{v_0}{\sigma_v^2}(\kappa - \rho\sigma_v\, iu - d)\frac{1 - e^{-d\tau}}{1 - \tilde{g}\,e^{-d\tau}}. \label{eq:B_heston}
\end{align}

The original Heston CF involves a complex-valued square root and a complex logarithm. Because the complex logarithm is multi-valued, numerical implementations must track the correct branch, which can lead to discontinuities in the CF as $u$ varies, a phenomenon known as the \textit{branch-cut problem}. This instability causes spurious oscillations in the numerically inverted PDF.

In \cite{heston:little_trap} Albrecher showed that a simple reformulation of Eq.~(\ref{eq:g_tilde}) eliminates the branch-cut issue entirely. The key insight is to define $\tilde{g}$ as the ratio $(\xi - d)/(\xi + d)$ where $\xi = \kappa - \rho\sigma_v\, iu$, and to enforce $\text{Re}(d) \geq 0$ in Eq.~(\ref{eq:d_heston}) by choosing the branch of the square root with non-negative real part. Under this convention, the logarithm in Eq.~(\ref{eq:A_heston}) remains on a consistent branch for all $u$, and the CF is smooth and numerically stable.\\




Given the CF $\varphi(u)$ of a continuous random variable $X$ with PDF $f(x)$, the \textit{Gil-Pelaez inversion theorem} \cite{heston:gil_pelaez} states that

\begin{equation}
    f(x) = \frac{1}{\pi}\int_0^\infty \text{Re}\!\left[\varphi(u)\,e^{-iux}\right]du.
    \label{eq:gilpelaez}
\end{equation}

This follows from the standard Fourier inversion theorem and the fact that $\varphi(-u) = \overline{\varphi(u)}$ for real-valued distributions, which halves the integration range. To evaluate Eq.~(\ref{eq:gilpelaez}) numerically, the integral is truncated at a finite upper limit $U = (N_\text{fft}-1)\eta$ and discretised on a uniform grid $u_j = j\eta$, $j = 0, 1, \ldots, N_\text{fft}-1$, with frequency step $\eta > 0$:

\begin{equation}
    f(x) \approx \frac{\eta}{\pi}\,\text{Re}\!\left[\sum_{j=0}^{N_\text{fft}-1} w_j\,\varphi(j\eta)\,e^{-ij\eta x}\right],
    \label{eq:gilpelaez_discrete}
\end{equation}

where $w_j$ are numerical quadrature weights.
The choice of quadrature weights $w_j$ in Eq.~(\ref{eq:gilpelaez_discrete}) determines the accuracy of the numerical integral. Simple rectangular (Euler) weights give first-order accuracy. The implementation instead uses \textit{Simpson's rule}, which achieves fourth-order accuracy by approximating the integrand as a piecewise quadratic polynomial. For $N_\text{fft}$ equally spaced points with step $\eta$, Simpson's rule assigns weights

\begin{equation}
    w_j = \frac{1}{3}\begin{cases} 1 & j = 0 \text{ or } j = N_\text{fft}-1, \\ 4 & j \text{ odd}, \\ 2 & j \text{ even}, \; j \neq 0, N_\text{fft}-1, \end{cases}
    \label{eq:simpson}
\end{equation}

The \textit{Discrete Fourier Transform} (DFT) of a sequence $\{x_j\}_{j=0}^{N-1}$ is defined as

\begin{equation}
    \hat{x}_k = \sum_{j=0}^{N-1} x_j\, e^{-i2\pi jk/N}, \quad k = 0, 1, \ldots, N-1.
    \label{eq:dft}
\end{equation}

The \textit{Fast Fourier Transform} (FFT) computes Eq.~(\ref{eq:dft}) in $\mathcal{O}(N\log N)$ operations rather than the $\mathcal{O}(N^2)$ required by direct summation. To apply the FFT to the inversion formula (\ref{eq:gilpelaez_discrete}), one evaluates Eq.~(\ref{eq:gilpelaez_discrete}) simultaneously at the grid of log-price values $x_k = b + k\lambda$, $k = 0, \ldots, N_\text{fft}-1$, for some lower bound $b$ and spacing $\lambda$. Inserting this into Eq.~(\ref{eq:gilpelaez_discrete}):

\begin{align}
    f(x_k) \approx \frac{\eta}{\pi}\,\mathrm{Re} \left[\sum_{j=0}^{N_{\mathrm{FFT}}-1}w_j\,\varphi(j\eta)\,e^{-ij\eta b}\,e^{-i2\pi\frac{\eta\lambda}{2\pi}jk}\right]
    \label{eq:fft_coupl}
\end{align}

Comparing with Eq.~(\ref{eq:dft}), the sum is a DFT if and only if the exponent is $e^{-i2\pi jk/N_\text{fft}}$, which requires

\begin{equation}
    \frac{\eta\lambda}{2\pi} = \frac{1}{N_\text{fft}} \;\Longrightarrow\; \lambda = \frac{2\pi}{N_\text{fft}\,\eta}.
    \label{eq:std_fft_constraint}
\end{equation}

This \textit{coupling constraint} links the frequency step $\eta$ and the spatial step $\lambda$: making the log-price grid finer (smaller $\lambda$) forces a larger $\eta$, which reduces the effective integration range in frequency and may introduce aliasing or truncation errors. This rigid coupling is the central limitation of the standard FFT approach to Fourier inversion.\\

The \textit{Fractional DFT} (frDFT) generalises Eq.~(\ref{eq:dft}) by replacing $1/N$ with an arbitrary real parameter $\alpha$:

\begin{equation}
    y_k = \sum_{j=0}^{N-1} x_j\, e^{-i2\pi\alpha jk}, \quad k = 0, \ldots, N-1.
    \label{eq:frdft}
\end{equation}

Comparing with Eq.~(\ref{eq:fft_coupl}), choosing $\alpha = \eta\lambda/(2\pi)$ evaluates the inversion integral at the spatial grid $\{x_k\}$ without any constraint between $\eta$ and $\lambda$. The user is therefore free to choose $\eta$ and $\lambda$ independently, completely decoupling frequency resolution from spatial resolution.
However, computing Eq.~(\ref{eq:frdft}) naively requires $\mathcal{O}(N^2)$ operations. 
Bluestein observed that the product $jk$ can be written as

\begin{equation}
    jk = \frac{j^2 + k^2 - (j-k)^2}{2},
    \label{eq:bluestein_id}
\end{equation}

which transforms Eq.~(\ref{eq:frdft}) into

\begin{equation}
    y_k = \underbrace{e^{-i\pi\alpha k^2}}_{c_k^*}\sum_j \underbrace{\bigl[x_j\, e^{-i\pi\alpha j^2}\bigr]}_{a_j = x_j c_j^*} \cdot \underbrace{e^{i\pi\alpha(j-k)^2}}_{h_{j-k}}= c_k^* \sum_j a_j\, h_{j-k}
    \label{eq:bluestein_conv}
\end{equation}

The sum over $j$ in Eq.~(\ref{eq:bluestein_conv}) is a \textit{linear convolution} of the sequence $a_j = x_j e^{-i\pi\alpha j^2}$ with the symmetric \textit{chirp kernel} $h_m = e^{i\pi\alpha m^2}$, evaluated at lags $m = j - k$. Linear convolution of length-$N$ sequences can be computed in $\mathcal{O}(N\log N)$ via the convolution theorem:
Note the FFT of $h$ is computed once so every $J$ trajectories
share one chirp and one kernel FFT, only the data-dependent part of the FFT 
operation differs per row.

\begin{equation}
    \text{Linear convolution: }\; a \star h = \mathcal{F}^{-1}\!\left[\mathcal{F}[a]\cdot\mathcal{F}[h]\right],
    \label{eq:conv_theorem}
\end{equation}

where the transforms are taken on a zero-padded array of length $M \geq 2N - 1$ (rounded up to the next power of two to allow the FFT) to prevent circular aliasing. The full frFFT algorithm is:

\begin{enumerate}
    \item Compute $a_j = x_j e^{-i\pi\alpha j^2}$ (phase-modulate input by conjugate chirp).
    \item Zero-pad $a$ to length $M = 2^{\lceil\log_2(2N-1)\rceil}$.
    \item Build the chirp kernel $h_{\text{pad}}$ of length $M$ in circular order: $h[0], h[1], \ldots, h[N-1]$ followed by zeros, then $h[N-1], \ldots, h[1]$.
    \item Compute $Y = \texttt{IFFT}(\texttt{FFT}(a_{\text{pad}}) \cdot \texttt{FFT}(h_{\text{pad}}))$.
    \item Multiply the first $N$ outputs: $y_k = e^{-i\pi\alpha k^2} Y_k$, $k = 0, \ldots, N-1$.
\end{enumerate}

\begin{algorithm}[H]
    \label{alg:frFTT}
\caption{Batched Fractional FFT via Bluestein's Convolution Identity}
\begin{algorithmic}[1]
\REQUIRE Batch $\mathbf{X}\in\mathbb{C}^{J\times N}$ (row $j$ = trajectory $j$'s frequency samples),
         fractional parameter $\alpha=\eta\lambda/(2\pi)$
\ENSURE  $\mathbf{Y}\in\mathbb{C}^{J\times N}$ with
         $Y_{j,k}=\sum_{n=0}^{N-1}X_{j,n}\,e^{-i2\pi\alpha nk}$,\; $k=0,\ldots,N-1,\; j=1,\ldots,J$

\LINECOMMENT{Identity: $nk=\tfrac12[n^2+k^2-(n-k)^2]$, so
  $e^{-i2\pi\alpha nk}=c_n^{*}c_k^{*}\,h_{n-k}$ with chirp $c_n=e^{i\pi\alpha n^2}$, $h_m=c_m$}

\LINECOMMENT{Step 0 --- precompute the chirp once, shared by every trajectory}
\STATE $c_n \leftarrow e^{i\pi\alpha n^2}$,\quad $n=0,\ldots,N-1$

\LINECOMMENT{Step 1 --- demodulate each row of the input by the conjugate chirp}
\FOR{$j=1$ \textbf{to} $J$}
    \STATE $a_{j,n}\leftarrow X_{j,n}\,c_n^{*}$,\quad $n=0,\ldots,N-1$
\ENDFOR

\LINECOMMENT{Step 2 --- padded length avoiding circular wrap-around}
\STATE $M\leftarrow 2^{\lceil\log_2(2N-1)\rceil}$

\LINECOMMENT{Step 3 --- zero-pad every row of $\mathbf{a}$ to length $M$}
\STATE $\mathbf{a}_{\mathrm{pad}}\leftarrow\mathbf{0}_{J\times M}$;\quad
       $\mathbf{a}_{\mathrm{pad}}[:,0{:}N]\leftarrow\mathbf{a}$

\LINECOMMENT{Step 4 --- build the kernel from the same chirp, laid out circularly
  ($h_{-m}=h_m$ places negative lags at the tail of length $M$)}
\STATE $\mathbf{h}_{\mathrm{pad}}\leftarrow\mathbf{0}_M$
\STATE $\mathbf{h}_{\mathrm{pad}}[0{:}N]\leftarrow[c_0,\ldots,c_{N-1}]$
       \COMMENT{lags $m=0,\ldots,N-1$}
\STATE $\mathbf{h}_{\mathrm{pad}}[M-N+1{:}M]\leftarrow[c_{N-1},\ldots,c_1]$
       \COMMENT{lags $m=-(N-1),\ldots,-1$}

\LINECOMMENT{Step 5 --- linear convolution via the convolution theorem;
  $\mathrm{FFT}(\mathbf{h}_{\mathrm{pad}})$ computed once and shared across all $J$ rows}
\STATE $\mathbf{conv}\leftarrow \mathrm{IFFT}\bigl(\mathrm{FFT}(\mathbf{a}_{\mathrm{pad}},\mathrm{axis}=1)\odot\mathrm{FFT}(\mathbf{h}_{\mathrm{pad}})\bigr)$

\LINECOMMENT{Step 6 --- demodulate by the chirp again, keep first $N$ entries}
\STATE $Y_{j,k}\leftarrow c_k^{*}\,\mathrm{conv}_{j,k}$,\quad $k=0,\ldots,N-1,\;j=1,\ldots,J$

\RETURN $\mathbf{Y}\in\mathbb{C}^{J\times N}$
\end{algorithmic}
\end{algorithm}


In the current implementation, the fractional parameter is set by default to $\alpha = 0.5/N_\text{fft}$, giving the spatial spacing

\begin{equation}
    \lambda = \frac{2\pi\alpha}{\eta} = \frac{4\pi}{N_\text{fft}\,\eta}.
    \label{eq:lambda_frfft}
\end{equation}

For $N_\text{fft} = 4096$ and $\eta = 0.25$, this yields $\lambda \approx 0.0031$, which is two orders of magnitude finer than the $\lambda_{\text{std}} = 2\pi/(N_\text{fft}\,\eta) \approx 0.0061$ that the standard FFT would enforce. The user can further tighten the grid by choosing a smaller $\alpha$ without any effect on the frequency resolution $\eta$.

% ─────────────────────────────────────────────────────────────────────────────
\subsection{The PDF Computation Pipeline}

In order to obtain the probability distribution function the code assembles the above components into a complete PDF computation. The following describes each stage of the pipeline.


\begin{itemize}
    \item \textbf{Grid Initialisation}. For each of the $J$ trajectories, a lower bound $b_j$ for the log-price evaluation grid is determined from an approximate mean and standard deviation of the forecast distribution. Since the exact moments depend on the full path of $v_t$, an approximation using the long-run variance $\theta$ is used:
    \begin{align}
        \text{approx\_mean}_j &= X_{0,j} + \left(\mu_j - \tfrac{1}{2}\theta_j\right)\tau, \label{eq:approx_mean}\\
        \text{approx\_std}_j  &= \sqrt{\max(\theta_j, 0)\,\tau}, \label{eq:approx_std}\\
        b_j &= \text{approx\_mean}_j - c\cdot\text{approx\_std}_j, \label{eq:lower_bound}
    \end{align}
    where $c = 5$ (the parameter \texttt{x\_width}) ensures that the grid spans approximately $\pm 5$ standard deviations around the forecast mean, capturing virtually all of the probability mass.
    \item \textbf{Phase Shift for the lower bound}. Equation (\ref{eq:gilpelaez_discrete}) evaluated at $x_k = b + k\lambda$ factors as
    \begin{equation}
    f(b + k\lambda) \approx \frac{\eta}{\pi}\,\text{Re}\!\left[\sum_{j=0}^{N_\text{fft}-1} \underbrace{\bigl(w_j\,\varphi(j\eta)\,e^{-ij\eta b}\bigr)}_{\text{frFFT input}} \cdot e^{-i2\pi\alpha jk}\right].
    \label{eq:phase_shifted}
    \end{equation}
    The factor $e^{-ij\eta b}$ shifts the output grid's origin to $b$ without changing the frequency-domain data. 
    \item \textbf{PDF extraction and clipping}. After the frFFT, the approximate PDF values are
    \begin{equation}
    f(x_k) \approx \frac{\eta}{\pi}\,\text{Re}\!\left[y_k\right], \quad k = 0, \ldots, N_\text{fft}-1.
    \label{eq:pdf_extract}
    \end{equation}
    Numerical inversion can produce small negative values due to discretisation error and finite truncation of the frequency integral. These artefacts are removed by clipping: \texttt{pdf\_fine = np.maximum(pdf\_fine, 0.0)}, which has negligible effect on well-resolved distributions.
    \item \textbf{Raw moments output}. When \texttt{n\_bins=None}, the method instead computes the mean and standard deviation directly from the fine frFFT grid using the rectangle rule:
    \begin{align}
    \mathbb{E}[X] &\approx \frac{\sum_k f(x_k)\,x_k\,\lambda}{\sum_k f(x_k)\,\lambda}, \label{eq:raw_mean}\\
    \text{Var}(X) &\approx \frac{\sum_k f(x_k)\,(x_k - \mathbb{E}[X])^2\,\lambda}{\sum_k f(x_k)\,\lambda}. \label{eq:raw_var}
    \end{align}
    The output is then a $(J, 2)$ array with columns containing the mean and standard deviation for each trajectory.
\end{itemize}




\begin{algorithm}[H]
    \label{alg:heston_pdf}
\caption{Heston PDF via Fractional FFT and Gil-Pelaez Inversion}
\begin{algorithmic}[1]
\REQUIRE Simulated terminal states $\{X_{T}^{(j)}, v_{T}^{(j)}\}_{j=1}^{J}$,
         Heston parameters $\{\mu^{(j)}, \kappa^{(j)}, \theta^{(j)},
         \sigma_v^{(j)}, \rho^{(j)}\}_{j=1}^{J}$
\REQUIRE Steps ahead $n_{\tau}$, time step $\Delta t$,
         FFT length $N$, frequency step $\eta$,
         fractional parameter $\alpha$, grid half-width $c$,
         number of output bins $K$ (or \texttt{None})
\ENSURE  Discretised conditional PDF or raw moments for each trajectory

\LINECOMMENT{Forecast horizon and x-domain spacing}
\STATE $\tau \leftarrow n_{\tau}\,\Delta t$
\STATE $\lambda \leftarrow 2\pi\alpha / \eta$

\LINECOMMENT{Per-trajectory lower bounds for the x-grid}
\FOR{$j = 1$ \textbf{to} $J$}
    \STATE $\tilde{\mu}^{(j)}   \leftarrow X_{T}^{(j)}
           + \bigl(\mu^{(j)} - \tfrac{1}{2}\theta^{(j)}\bigr)\tau$
           \COMMENT{approximate centre}
    \STATE $\tilde{\sigma}^{(j)} \leftarrow \sqrt{\theta^{(j)}\,\tau}$
           \COMMENT{approximate spread}
    \STATE $b^{(j)} \leftarrow \tilde{\mu}^{(j)} - c\,\tilde{\sigma}^{(j)}$
           \COMMENT{grid lower bound}
\ENDFOR

\LINECOMMENT{Frequency grid and Simpson quadrature weights}
\STATE $u_n \leftarrow n\,\eta$,\quad $n = 0, \ldots, N-1$
\STATE $w_0 \leftarrow 1/3$;\quad
       $w_n \leftarrow 4/3$ for $n$ odd;\quad
       $w_n \leftarrow 2/3$ for $n$ even, $n \notin \{0, N-1\}$;\quad
       $w_{N-1} \leftarrow 1/3$

\LINECOMMENT{Evaluate Heston characteristic function (eq. \ref{eq:cf_heston}) for all trajectories}
\STATE $\boldsymbol{\varphi} \leftarrow
       \textsc{HestonCF}\!\left(
         \mathbf{u},\,\tau,\,
         \{X_{T}^{(j)}\},\{v_{T}^{(j)}\},
         \{\mu^{(j)}\},\{\kappa^{(j)}\},\{\theta^{(j)}\},\{\sigma_v^{(j)}\},\{\rho^{(j)}\}
       \right)
       \in \mathbb{C}^{J \times N}$

\LINECOMMENT{Phase-shift to lower bound $b^{(j)}$, then weight by quadrature}
\STATE $\mathbf{Z}^{(j)}_n \leftarrow
       w_n\,\varphi^{(j)}_n\,
       e^{-i\,u_n\,b^{(j)}}$,\quad
       $j = 1,\ldots,J,\quad n = 0,\ldots,N-1$

\LINECOMMENT{Batched frFFT inversion (Algorithm~\ref{alg:frFTT})}
\STATE $\mathbf{Y} \leftarrow
       \textsc{FrFFT\_Batch}\!\left(\mathbf{Z},\,\alpha\right)
       \in \mathbb{C}^{J \times N}$

\LINECOMMENT{Recover density values and clip numerical artefacts}
\STATE $f^{(j)}_k \leftarrow
       \max\!\left(0,\;
         \dfrac{\eta}{\pi}\,\mathrm{Re}\!\left[Y^{(j)}_k\right]
       \right)$,\quad
       $k = 0,\ldots,N-1$

\LINECOMMENT{x-grid for each trajectory}
\STATE $x^{(j)}_k \leftarrow b^{(j)} + k\,\lambda$,\quad
       $k = 0,\ldots,N-1$

\IF{$K$ is \texttt{None}}
    \LINECOMMENT{Return raw moments via numerical integration}
    \FOR{$j = 1$ \textbf{to} $J$}
        \STATE $C^{(j)} \leftarrow \lambda\sum_{k} f^{(j)}_k$
               \COMMENT{normalisation constant}
        \STATE $\bar{x}^{(j)} \leftarrow
               \lambda\sum_{k} x^{(j)}_k\,f^{(j)}_k \;/\; C^{(j)}$
        \STATE $s^{(j)} \leftarrow
               \sqrt{\lambda\sum_{k}
                 \bigl(x^{(j)}_k - \bar{x}^{(j)}\bigr)^{2}
                 f^{(j)}_k \;/\; C^{(j)}}$
    \ENDFOR
    \RETURN $\bigl[\bar{x}^{(j)},\,s^{(j)}\bigr]_{j=1}^{J}$
\ENDIF

\RETURN $\mathbf{p} = f^{(j)} * x^{(j)} \in \mathbb{R}^{J \times K}$

\end{algorithmic}
\end{algorithm}

Moreover, the code provides a reference implementation based on direct Monte Carlo simulation. For each of the $J$ starting conditions $(X_{0,j}, v_{0,j})$, a large number \texttt{mc\_sims} of independent forward paths are generated using the full-truncation Euler--Maruyama scheme over \texttt{n\_steps} steps. The terminal log-prices are then histogrammed onto the shared bin structure to produce an empirical approximation of the PDF.

The Monte Carlo estimator is consistent (converges to the true PDF as \texttt{mc\_sims} $\to\infty$) but slow: its computational cost is $\mathcal{O}(J \cdot \texttt{mc\_sims} \cdot \texttt{n\_steps})$. The frFFT method, by contrast, evaluates the exact Heston CF and inverts it in $\mathcal{O}(J \cdot N_\text{fft}\log N_\text{fft})$ operations, making it orders of magnitude faster for large $J$.

\begin{algorithm}[H]
\caption{Monte-Carlo Conditional Heston PDF}
\label{alg:heston_mc_pdf}
\begin{algorithmic}[1]

\REQUIRE Current states $\{X_0^{(j)},v_0^{(j)}\}_{j=1}^{J}$,
Heston parameters
$\{\kappa^{(j)},\theta^{(j)},\sigma_v^{(j)},\rho^{(j)},\mu^{(j)}\}_{j=1}^{J}$,
forecast horizon $\tau=n_\tau\Delta t$,
Monte-Carlo sample size $M$

\ENSURE Terminal log-price samples
$\mathbf{X}\in\mathbb{R}^{J\times M}$

\FOR{$j=1$ \textbf{to} $J$}

    \STATE Initialise a CIR process from $v_0^{(j)}$ using
    $(\kappa^{(j)},\theta^{(j)},\sigma_v^{(j)})$.

    \STATE Evolve the variance process
    (\ref{eq:var_sde_heston}) over $[0,\tau]$
    using Andersen's Quadratic--Exponential (QE) scheme,
    simultaneously computing
    \[
        v_\tau,
        \qquad
        \int_0^\tau v(s)\,ds
    \]
    for all $M$ Monte-Carlo paths.

\STATE   Compute the conditional log-price increment
\[
Y=
-\frac12 I_v
+\frac{\rho}{\sigma_v}
\Bigl(v_\tau-v_0-\kappa(\theta\tau-I_v)\Bigr)
+\sqrt{(1-\rho^2)I_v}\,Z,
\]
where $I_v=\int_0^\tau v(s)\,ds$.

\STATE Set
\[
X_\tau=X_0+Y+\mu\tau.
\]

    \STATE Store the $M$ simulated terminal values as row $j$ of
    $\mathbf{X}$.

\ENDFOR

\RETURN $\mathbf{X}$

\end{algorithmic}
\end{algorithm}

So the \texttt{HestonSimulator} integrates the following chain of mathematical techniques into a single coherent pipeline. 
The Heston SDE system (\ref{eq:logprice_sde_heston})--(\ref{eq:correlation_heston}) is discretised using correlated Brownian increments constructed via Cholesky decomposition, 
with the full-truncation Milstein scheme for the variance and Euler--Maruyama for the log-price. 
The conditional PDF of future log-prices is obtained by evaluating the closed-form Heston CF in its numerically 
stable formulation and then applying the Gil-Pelaez inversion formula (\ref{eq:gilpelaez_discrete}) with Simpson quadrature weights.
The resulting discrete sum is evaluated efficiently across all $N_\text{fft}$ spatial grid points simultaneously by expressing it as a fractional DFT and computing
it via Bluestein's chirp convolution algorithm, which reduces the cost from $\mathcal{O}(N_\text{fft}^2)$ to $\mathcal{O}(N_\text{fft}\log N_\text{fft})$ 
while completely decoupling the frequency resolution $\eta$ from the spatial resolution $\lambda$.\\
The Monte Carlo estimator is validated against the frFFT method over $10000$ sample trajectories obtaining a $100\%$ of acceptance rate of the Kolmogorov-Smirnov test at a $5\%$ significance level.
Below are presented four samples for visual inspection of the two methods, showing a very good agreement between them.

\begin{figure}[H]
    \centering
    \includegraphics[width=\textwidth]{Heston/images/mc_heston_pdf.png}
    \caption{Comparison of the Heston PDF obtained via the frFFT method (red line) and the Monte Carlo estimator (blue histogram) for four sample trajectories. }
    \label{fig:heston_mc_plots}
\end{figure}



% ─────────────────────────────────────────────────────────────────────────────
% Section: Numerical Validation of the Fractional FFT via the
%          Black-Scholes Characteristic Function
% ─────────────────────────────────────────────────────────────────────────────

\section{Numerical Validation of the Fractional FFT via the
         Black-Scholes Characteristic Function}
\label{sec:frfft_validation}

The implementation of the fractional Fast Fourier Transform presented in the
preceding section constitutes a complex numerical pipeline: it involves a
quadrature approximation of an improper integral, a chirp-modulated linear
convolution evaluated through nested FFT calls, and a subsequent restriction
of the recovered density to a bounded support.  Each step introduces
approximation error, and it is therefore essential to verify that the full
pipeline produces numerically correct results before it is deployed inside the
Heston simulator.

The validation strategy adopted here exploits a key property of the
Black-Scholes model: under Geometric Brownian Motion the conditional
log-price distribution has a \emph{known closed form}.  Specifically, given
the log-price $X_t$ at time $t$, the one-step-ahead log-price
$X_{t+\tau} \mid X_t$ follows a Gaussian distribution with analytically
available parameters.  By feeding the Black-Scholes characteristic function
into the fractional FFT inversion routine and comparing the recovered density
to that closed-form Gaussian, one obtains a ground-truth benchmark that is
model-free with respect to the numerical algorithm under test.  The
experiment is therefore not a Monte Carlo approximation: it is an exact
algebraic comparison, and any discrepancy between the two densities is
attributable entirely to numerical error in the frFFT machinery.\\


Under the Black-Scholes model the log-price evolves according to the
stochastic differential equation
\begin{equation}
    dX_t \;=\; \Bigl(\mu - \tfrac{1}{2}\sigma^2\Bigr)\,dt \;+\; \sigma\,dW_t,
    \label{eq:bs_sde}
\end{equation}
where $\mu$ is the drift, $\sigma > 0$ is the volatility, and $W_t$ is a
standard Brownian motion.  Over a fixed horizon $\tau > 0$ the increment
$X_{t+\tau} - X_t$ is normally distributed with mean
$\bigl(\mu - \tfrac{1}{2}\sigma^2\bigr)\tau$ and variance $\sigma^2\tau$,
so the conditional distribution reads
\begin{equation}
    X_{t+\tau} \mid X_t \;\sim\;
    \mathcal{N}\!\bigl(\mu_{\mathrm{BS}},\,\sigma_{\mathrm{BS}}^2\bigr),
    \qquad
    \mu_{\mathrm{BS}} = X_t + \bigl(\mu - \tfrac{1}{2}\sigma^2\bigr)\tau,
    \quad
    \sigma_{\mathrm{BS}}^2 = \sigma^2\tau.
    \label{eq:bs_normal}
\end{equation}
The corresponding characteristic function (CF) is
\begin{equation}
    \varphi_{\mathrm{BS}}(u)
    \;=\;
    \mathbb{E}\bigl[e^{iuX_{t+\tau}}\bigr]
    \;=\;
    \exp\!\Bigl(iu\,\mu_{\mathrm{BS}} \;-\; \tfrac{1}{2}\sigma_{\mathrm{BS}}^2\,u^2\Bigr),
    \qquad u \in \mathbb{R}.
    \label{eq:bs_cf}
\end{equation}
This is the CF of a Gaussian and is entire in $u$, which makes the
Gil-Pelaez inversion numerically well-conditioned: the integrand decays
super-exponentially as $|u| \to \infty$, so the truncated quadrature sum
converges rapidly.

\subsection{Validation Design}
\label{subsec:validation_design}

The experiment is structured in three sequential validation steps, each
targeting a different layer of the computational pipeline.\\


A set of $J = 10000$ independent trajectories is simulated using the 
same algorithm described in section \ref{sec:bs_experiments}, in particular \ref{alg:y condition}, with $N = 252$ time steps over a horizon of
$T = 1$ year, giving a time-step size of $dt = T/N = 0.004$ years.  The
per-trajectory drift and volatility parameters are drawn uniformly:
$\mu \sim \mathcal{U}(0,\, 0.1)$ and $\sigma \sim \mathcal{U}(0.1,\, 1.2)$,
while the initial log-price satisfies $X_0 \sim \mathcal{U}(0,\, 1)$.  With
a one-step-ahead forecast horizon $\tau = dt = 0.004$ years.\\

The inversion is performed with $N_{\mathrm{FFT}} = 4096$ frequency nodes
and a frequency-domain step $\eta = 0.25$.  Following algorithm \ref{alg:frFTT}. The fractional parameter is set to
$\alpha = 2/N_{\mathrm{FFT}} \approx 4.88 \times 10^{-4}$, which yields an
$x$-domain spacing of
\begin{equation}
    \lambda \;=\; \frac{2\pi\alpha}{\eta} \;\approx\; 0.01227.
    \label{eq:lambda}
\end{equation}
The lower bound of the $x$-grid for each trajectory $k$ is placed at
$b_k = \mu_{\mathrm{BS},k} - 5\,\sigma_{\mathrm{BS},k}$, centring the
grid on the Gaussian peak and ensuring that $\pm 5\sigma$ tails are always
covered.\\

\paragraph{Step A: Analytical moments sanity check.}
Before engaging the frFFT at all, the first validation step confirms that the
method used to obtain the analytical probability distribution function correctly implements
equation~\eqref{eq:bs_normal}.  The
method returns a $(J \times 2)$ array whose columns are the analytical mean
$\mu_{\mathrm{BS},k}$ and standard deviation $\sigma_{\mathrm{BS},k}$ for
each trajectory $k$.  These are compared against the values computed directly
from the simulation parameters and the closed-form expressions in
equation~\eqref{eq:bs_normal}.

The results, summarised in Table~\ref{tab:step_a}, show that the analytical
moments returned by the code implementation are \emph{bit-for-bit identical} to the
directly computed values: the absolute difference in both mean and standard
deviation is exactly zero across all four trajectories.  This is expected,
since both quantities are computed from the same floating-point operations on
the same arrays; nevertheless it confirms that there is no implementation
error in the formula for the analytical parameters.

\begin{table}[H]
\centering
\caption{Step A: comparison between directly computed analytical parameters
and those returned by code simulation for the first 4 trajectories.
All differences are zero to double-precision.}
\label{tab:step_a}
\begin{tabular}{ccccccc}
\toprule
$k$ & $\mu_{\mathrm{BS},k}$ & \texttt{computed} mean & $|\Delta\mu|$
    & $\sigma_{\mathrm{BS},k}$ & \texttt{computed} std & $|\Delta\sigma|$ \\
\midrule
0 & 0.8533110 & 0.8533110 & $0$ & 0.0232011 & 0.0232011 & $0$ \\
1 & 0.1042671 & 0.1042671 & $0$ & 0.0459892 & 0.0459892 & $0$ \\
2 & 0.6879010 & 0.6879010 & $0$ & 0.0403615 & 0.0403615 & $0$ \\
3 & 0.7827374 & 0.7827374 & $0$ & 0.0796743 & 0.0796743 & $0$ \\
\bottomrule
\end{tabular}
\end{table}

\paragraph{Step B: frFFT Inversion.}
The core validation step applies the Gil-Pelaez inversion formula to the
Black-Scholes CF of equation~\eqref{eq:bs_cf}.  The inversion approximates
the probability density function as
\begin{equation}
    f(x_k) \;\approx\;
    \frac{\eta}{\pi}\,
    \mathrm{Re}\!\left[
        \sum_{j=0}^{N_{\mathrm{FFT}}-1}
        w_j\,\varphi_{\mathrm{BS}}(j\eta)\,e^{-ij\eta x_k}
    \right],
    \qquad x_k = b + k\lambda,
    \label{eq:gp_discrete}
\end{equation}
where $w_j$ are Simpson quadrature weights.  The sum over $j$ is evaluated
simultaneously at all $k = 0, \ldots, N_{\mathrm{FFT}} - 1$ grid points
using the batched fractional FFT described in algorithm \ref{alg:frFTT},
which implements Bluestein's identity.  Concretely, defining the frequency-domain
input
\begin{equation}
    \tilde{x}_j \;=\; w_j\,\varphi_{\mathrm{BS}}(j\eta)\,e^{-ij\eta b},
    \label{eq:fft_input}
\end{equation}
equation~\eqref{eq:gp_discrete} reduces to the fractional DFT
\begin{equation}
    Y[k] \;=\;
    \sum_{j=0}^{N-1} \tilde{x}_j\,e^{-i2\pi\alpha jk},
    \qquad
    f(x_k) = \frac{\eta}{\pi}\,\mathrm{Re}[Y[k]],
    \label{eq:frfft_sum}
\end{equation}

An important numerical subtlety arises from the choice $\alpha = 2/N$.
In this regime, the frFFT evaluates the standard DFT at even-indexed
frequencies: $Y[k] = \hat{X}[2k]$, where $\hat{X}$ denotes the length-$N$
DFT of the input sequence.  A consequence of this sub-sampling in frequency
space is that the second half of the output grid aliases onto the first half,
effectively doubling the contribution of the Gaussian peak.  As a result, the
raw integral
\begin{equation}
    \int_{-\infty}^{\infty} f(x)\,dx
    \;\approx\;
    \lambda \sum_{k=0}^{N-1} f(x_k)
    \;\approx\; 2,
    \label{eq:raw_integral}
\end{equation}
instead of the theoretically correct value of one.  The numerical output
confirms this precisely: the raw integral equals $2.0000$ for every
trajectory, as shown below.

\begin{center}
\texttt{t0:2.0000 \quad t1:2.0000 \quad t2:2.0000 \quad t3:2.0000}
\end{center}

This factor of two is \emph{systematic and entirely predictable}: it does
not represent an error in the inversion formula, but rather a known property
of the fractional parameter choice $\alpha = 2/N$.  The \texttt{HestonSimulator}
handles it transparently by normalising the recovered density after
restricting it to the relevant support, exactly as one would normalise a
histogram.  Specifically, the density is zeroed outside the interval
$[\mu_{\mathrm{BS}} - 5\sigma_{\mathrm{BS}},\,
  \mu_{\mathrm{BS}} + 5\sigma_{\mathrm{BS}}]$
and then rescaled so that $\lambda \sum_k f(x_k) = 1$.  This restriction also
has a direct physical interpretation: a $\pm 5\sigma$ interval captures
$99.9999\%$ of the Gaussian probability mass, so no meaningful probability is
discarded.\\

Following normalisation, the mean and standard deviation of the recovered
density are computed numerically from the frFFT grid as
\begin{equation}
    \hat{\mu}_k = \lambda\sum_{j \in \mathcal{S}_k} f(x_j)\,x_j,
    \qquad
    \hat{\sigma}_k = \sqrt{\lambda\sum_{j \in \mathcal{S}_k}
                           f(x_j)\,(x_j - \hat{\mu}_k)^2},
    \label{eq:frfft_moments}
\end{equation}
where $\mathcal{S}_k$ is the index set of grid points within the support.
Table~\ref{tab:step_b} reports these moments alongside the analytical
parameters and the corresponding absolute errors.

\begin{table}[H]
\centering
\caption{Step B: moment accuracy of the frFFT-recovered density for the first 4 trajectories.  The
frFFT mean matches the analytical mean to within $\approx 10^{-7}$ and the
standard deviation to within $\approx 10^{-7}$.}
\label{tab:step_b}
\begin{tabular}{ccccccc}
\toprule
$k$ & $\hat{\mu}_k$ & $\mu_{\mathrm{BS},k}$ & $|\Delta\mu|$
    & $\hat{\sigma}_k$ & $\sigma_{\mathrm{BS},k}$ & $|\Delta\sigma|$ \\
\midrule
0 & 0.8533109 & 0.8533110 & $7.02 \times 10^{-8}$ & 0.0232009 & 0.0232011 & $2.05 \times 10^{-7}$ \\
1 & 0.1042671 & 0.1042671 & $2.94 \times 10^{-8}$ & 0.0459890 & 0.0459892 & $2.38 \times 10^{-7}$ \\
2 & 0.6879009 & 0.6879010 & $7.36 \times 10^{-8}$ & 0.0403612 & 0.0403615 & $3.08 \times 10^{-7}$ \\
3 & 0.7827373 & 0.7827374 & $8.18 \times 10^{-8}$ & 0.0796737 & 0.0796743 & $5.95 \times 10^{-7}$ \\
\bottomrule
\end{tabular}
\end{table}

The absolute error in the mean is at most $8.2 \times 10^{-8}$, while the
error in the standard deviation does not exceed $6.0 \times 10^{-7}$.  Both
figures are several orders of magnitude smaller than $\lambda \approx 0.012$,
the grid spacing of the frFFT output, which represents the fundamental
resolution limit of the discretisation.  The moment errors are therefore
entirely consistent with the quadrature discretisation error of the
Gil-Pelaez inversion and contain no evidence of a systematic implementation
defect.

\paragraph{Step C: Distribution Metrics.}
\label{subsec:step_c}

The final validation step compares the full shape of the frFFT-recovered
density with the analytical Gaussian at every active grid point within the
support.  Because the frFFT grid spacing $\lambda \approx 0.012$ is
considerably coarser than the Gaussian standard deviations
$\sigma_{\mathrm{BS}} \in [0.023, 0.080]$, there are on average only
$2 \times 5\sigma_{\mathrm{BS}} / \lambda \approx 38$ active grid points per
trajectory within the $\pm 5\sigma$ support.  Re-binning onto a finer histogram
would leave most bins empty, artificially inflating divergence metrics such as
Jensen-Shannon distance.  The correct approach is therefore to treat the two
densities as discrete probability vectors evaluated at the \emph{same}
$38$ quadrature points, converting density values to probability masses via
$p_j = f(x_j) \cdot \lambda$.

Four distribution metrics are computed: the Kolmogorov-Smirnov (KS)
statistic $D = \max_j |F_{\mathrm{frFFT}}(x_j) - F_{\mathrm{Gauss}}(x_j)|$,
where $F$ denotes the discrete CDF; the Jensen-Shannon distance
$\mathrm{JS}(p \,\|\, q) = \sqrt{D_{\mathrm{KL}}(p \,\|\, m) / 2
+ D_{\mathrm{KL}}(q \,\|\, m)/2}$ with $m = (p+q)/2$; the total variation
distance $\mathrm{TV} = \tfrac{1}{2}\sum_j |p_j - q_j|$; and the
Hellinger distance $H = \sqrt{\tfrac{1}{2}\sum_j (\sqrt{p_j} - \sqrt{q_j})^2}$.\\
The average values of the distribution metrics between the frFFT density and the
analytical Gaussian for the $10,000$ trajectories are in the order of $10^{-8}$ for the Jensen-Shannon, TV distance and Hellinger distance
while the average value of the KS test's pvalues is around $1.0$ suggesting no relevant difference between the 
analytical distribution and the one obtained from the fractional fft transform.\\

This result has a precise
interpretation: at each of the 38 active quadrature points, the probability
mass $p_j = f_{\mathrm{frFFT}}(x_j) \cdot \lambda$ and
$q_j = f_{\mathrm{Gauss}}(x_j) \cdot \lambda$ agree to within
double-precision floating-point resolution, so no measurable distance exists
between the two discrete distributions.\\

To probe the numerical error more finely the plots of the probability densities obtained from the frFFT and the analytical Gaussian are shown in Figure~\ref{fig:frfft_plots}. 
Moreover, Figure~\ref{fig:frfft_residuals}
plots the pointwise residual
$r(x_k) = f_{\mathrm{frFFT}}(x_k) - f_{\mathrm{Gauss}}(x_k)$
at every active grid point in the support for four randomly sampled trajectories.

\begin{figure}[H]
    \centering
    \includegraphics[width=\textwidth]{Heston/images/fft_plots.png}
    \caption{Comparison of the frFFT-recovered density (red dashed) and the analytical Gaussian (blue) for four randomly sampled trajectories.}
    \label{fig:frfft_plots}
\end{figure}


\begin{figure}[H]
    \centering
    \includegraphics[width=\textwidth]{Heston/images/fft_experiment.png}
    \caption{Pointwise residuals between the normalised frFFT density and the
    analytical Gaussian density, evaluated at the active fine-grid points
    within $\pm 5\sigma$ of each trajectory's mean.  The vertical axes carry
    exponents of $1^{-15}$, confirming that the residuals are at the level of double-precision
    machine epsilon.}
    \label{fig:frfft_residuals}
\end{figure}

The figure reveals several important features.  First, the magnitude of the
residuals is around $1^{-15}$, which lies within one to
two orders of magnitude of the double-precision machine epsilon
$\varepsilon_{\mathrm{mach}} = 2^{-52} \approx 2.22 \times 10^{-16}$. In other words,
the frFFT recovers the Gaussian density \emph{to the last representable bit}.

Second, the residuals are visibly oscillatory and change sign multiple times
across the support, with no discernible systematic trend or bias.  This
sign-changing, zero-mean pattern is the hallmark of floating-point rounding
error, which is essentially random in sign and magnitude, as opposed to
truncation error or quadrature bias, which would manifest as a residual that
is consistently positive or negative, or that grows monotonically toward the
boundaries.  The complete absence of such structured deviations confirms that
the Simpson quadrature weights and the frequency truncation at $j\eta_{\max}
= N_{\mathrm{FFT}} \cdot \eta = 1024$ are adequate for the narrow Gaussian
distributions encountered here: the CF has already decayed to
$|\varphi_{\mathrm{BS}}(1024)| = \exp(-\tfrac{1}{2}\sigma_{\mathrm{BS}}^2
\cdot 1024^2) \approx 10^{-2400}$, making the truncation error entirely
negligible.


\subsection{Summary and Interpretation}
\label{subsec:validation_summary}

The three-step experiment establishes the correctness of the frFFT
implementation with high confidence.  The main findings can be stated as
follows.

\begin{enumerate}
    \item \textbf{Moment accuracy}: the frFFT recovers the mean and standard
    deviation of the Gaussian to within $10^{-7}$, which is well below the
    grid-spacing resolution $\lambda \approx 0.012$, demonstrating that the
    numerical quadrature of the Gil-Pelaez integral is essentially exact.

    \item \textbf{Distributional identity}: all distribution metrics (KS, TV,
    Hellinger) are identically zero at the level of double-precision
    arithmetic, confirming that the recovered discrete probability masses
    agree point-by-point with those of the analytical Gaussian.

    \item \textbf{Machine-epsilon residuals}: the pointwise density residual
    lies between $10^{-15}$ and $10^{-14}$, consistent with the rounding
    error floor of machine precision, with no systematic bias.

    \item \textbf{Aliasing is benign}: the factor-of-two inflation of the raw
    integral, which is a predictable consequence of the choice
    $\alpha = 2/N_{\mathrm{FFT}}$, is fully corrected by the
    normalisation step and has no effect on the final PDF.
\end{enumerate}

These results validate the frFFT inversion routine for use in the Heston
model, where the characteristic function is considerably more complex but the
numerical pipeline is identical.  Since the Black-Scholes CF is an extremely
well-conditioned test case -- it is entire, super-exponentially decaying, and
possesses an exact closed-form inverse -- the experiment places an upper bound
on the algorithmic error: any additional inversion error encountered in the
Heston case will arise from the slower decay of the Heston CF and from the
presence of the stochastic variance term, not from a defect in the frFFT
machinery itself.





















\section{Scoring Rule ForGAN results}

The experiment has been conducted using the Heston trajectories and the frFFT-based PDF estimator described in \ref{sec:frfft_validation}.\\
The training procedure implemented is designed to run on a high-performance computing (HPC) of the CINECA nodes. In particular, the training script is executed on a node equipped with \textbf{8 CPU cores and 1 GPU}, as available on the Leonardo Booster partition as explained in the setup discussion \ref{setup}.\\
The script follows a fully automated pipeline that (i) generates a large synthetic Heston dataset in chunks, (ii) stores it in a binary streaming format, (iii) builds a memory-mapped dataset to avoid loading the full sample set into RAM, (iv) trains the SRForGAN generator on the GPU, and (v) saves the learned weights together with the full training configuration for reproducibility.\\

The code implements the following workflow:
\begin{enumerate}
    \item Generate synthetic Heston trajectories in chunks by using equation described in \ref{eq:milstein_full_trajectory}.
    \item Store each chunk in a flat binary file with a JSON header describing the record size.
    \item Build a memory-mapped dataset that reads trajectories lazily from disk.
    \item Split the dataset into training and validation subsets.
    \item Instantiate the SRForGAN model and select an LSTM generator.
    \item Train the model on the GPU using the Energy Score objective as defined in algorithm \ref{alg:srforgan_train}.
    \item Save the trained weights and configuration in the model directory.
    \item Test the model on unseen data and evaluate performance metrics.
\end{enumerate}

\subsection{Steps for experiment reproducibility}
The first stage constructs a large synthetic dataset based on the Heston stochastic-volatility model. 
The code simulates \text{10,000,000} trajectories, each of length \text{252 + 1}. To avoid memory overflow, the trajectories are generated in chunks of size $100,000$.

Each chunk is produced by the Heston simulator described in previous section \ref{eq:milstein_full_trajectory} and appended to a binary file.
A companion JSON file stores:
\begin{itemize}
    \item the record length, equal to \texttt{N\_STEPS = 252 + 1} where the 252 values are used as the conditioning window and the last value is used as the prediction target,
    \item the number of stored trajectories \text{J = 100,000}.
\end{itemize}

This structure allows the file to be accessed sequentially or through memory mapping without loading the entire dataset into memory.
The following parameters have been used to generate Heston trajectories:
\begin{table}[H]
    \label{tab:heston_training_params}
\centering
\begin{tabular}{lc}
\hline
Parameter & Value \\
\hline
$x_0$      & 0.0 \\
$\mu$      & 0.0 \\
$v_0$      & 0.5 \\
$\kappa$   & 2.0 \\
$\theta$   & 0.5 \\
$\sigma_v$ & 0.9 \\
$\rho$     & -0.7 \\
\texttt{Condition\_trajectory\_length} & 252 \\
\texttt{Forecast\_horizon}     & 10 \\
\hline
\end{tabular}
\caption{Heston parameters}
\end{table}
Where $x_0$ is the initial value of the log-price process, $\mu$ is the drift, $v_0$ is the initial variance, $\kappa$ is the mean-reversion speed of the variance, $\theta$ is the long-term mean of the variance, $\sigma_v$ is the volatility of the variance process, and $\rho$ is the correlation between the Brownian motions driving the log-price and variance.
From the parameters values, the Feller condition is satisfied (in fact, $\kappa \theta = 2.0 \times 0.5 = 1.0 > 0.9 = \sigma_v^2$), ensuring that the variance process remains positive. 


\begin{algorithm}[H]
\caption{Chunked generation of the Heston dataset}
\label{alg:heston_experiment}
\begin{algorithmic}[1]
\REQUIRE parameter ranges, time horizon $T$, number of steps $N$, total number of trajectories $J_{total}$, chunk size $J$
\STATE create \texttt{data/} and \texttt{models/} directories if needed
\STATE initialize binary file and JSON header
\FOR{$i = 1$ to $J_{total}/J$}
    \STATE simulate $J$ Heston trajectories with fixed seed $42+i$
    \STATE convert trajectories to 32-bit floating point format
    \STATE append the raw bytes to the binary file
    \STATE update the number of stored records in the JSON header
\ENDFOR
\end{algorithmic}
\end{algorithm}

A memory-mapped dataset is constructed from the binary file. Each stored record corresponds to a full trajectory
\[
[X_0, X_1, \dots, X_N],
\]
where the first $N$ values are used as the conditioning window and the last value is used as the prediction target. In other words, the dataset returns:
\[
c = [X_0, X_1, \dots, X_{N-10}], \qquad y = X_N.
\]


\begin{algorithm}[H]
\caption{Lazy access to the memory-mapped dataset}
\label{alg:memmap_dataset}
\begin{algorithmic}[1]
\REQUIRE binary trajectory file, JSON header with record length and number of records
\STATE read metadata from the JSON header
\STATE store the file path and defer memmap creation
\STATE \textbf{function} getitem(index)
\IF{memmap has not been opened yet}
    \STATE open the binary file as a NumPy memmap
\ENDIF
\STATE read trajectory $[X_0, \dots, X_N]$ at the requested index
\STATE set condition $c \leftarrow [X_0, \dots, X_{N-10}]$
\STATE set target $y \leftarrow X_N$
\STATE \textbf{return} $(y,c)$
\end{algorithmic}
\end{algorithm}


Once the dataset has been created, it is split into training and validation subsets using $80\%$ of the dataset for training and the remaining $20\%$ for validation.
The split is randomized with a fixed seed to preserve reproducibility.\\

The implemented SRForGAN model \ref{alg:srforgan_train} is initialized using the following configuration:
\begin{itemize}
    \item latent noise dimension \texttt{Z\_NOISE\_DIM = 64},
    \item scoring rule \texttt{energy} \ref{alg:energy_score},
    \item early stopping patience $ 10$ epochs,
    \item batch size  $128$,
    \item learning rate \texttt{$LR = 10^{-3}$}.
\end{itemize}

Moreover, given the complexity of the Heston process with respect to the Black-Scholes case, 
the generator architecture is designed to capture temporal dependencies in the conditioning window
by using the LSTM variant of the generator. In this variant the generator first extract hidden states from the conditioning window through an LSTM layer and then concatenates the final hidden state with 
the latent noise vector which will become the input of a feedforward network to produce the output prediction.
The hidden dimension of the LSTM layer is set to 64.\\

\begin{table}[H]
\centering
\caption{Generator Architecture and Hyperparameters Summary}
\label{tab:srforgan_architecture}
\begin{tabular}{lll}
\hline
\textbf{Parameter}         & \textbf{Generator}                  & \textbf{Critic (Discriminator)} \\ \hline
Input Dimension            & 64 (Latent) + 252 (Cond.)           & --                              \\
LSTM Layers                & 128                                 & --                              \\
Feedforward Layers         & 128+64                                 & --                              \\
Output Dimension           & 1                                  & --                              \\
Activation Function        & Leaky ReLU                         & --                              \\
Dropout                    & 0.0                                & --                              \\
\multicolumn{3}{c}{\textbf{Global Training Settings}}                                      \\ \hline
Batch Size                 & \multicolumn{2}{l}{32}                                            \\
max\_epochs                & \multicolumn{2}{l}{100}                                          \\
early\_stopping\_patience & \multicolumn{2}{l}{10}                                           \\
learning\_rate            & \multicolumn{2}{l}{$10^{-3}$}                           \\         \\
\hline
\end{tabular}
\end{table}



\subsection{Test results}
The trained model is evaluated on a test set given by \text{100,000} trajectories that were not seen during training by using the same Heston parameters adopted for the training process and described in table \ref{tab:heston_training_params}. 
The evaluation metrics are computed by comparing the distribution of the generated predictions with the distribution of the true targets in the test set, using the Heston Monte Carlo simulations (and the frFFT-based PDF estimator for the visualization) described in \ref{sec:frfft_validation}.\\

\begin{table}[h]
\centering
\begin{tabular}{lc}
\hline
Metric & Value \\
\hline
KS test acceptance rate & 87.500\% \\
Jensen--Shannon distance & 0.121850 \\
Hellinger distance & 0.104270 \\
Total variation distance & 0.102780 \\
Earth mover's distance & 0.000707 \\
\hline
\end{tabular}
\caption{Distribution comparison metrics}
\end{table}




\begin{figure}[H]
    \centering
    \includegraphics[width=1.0\textwidth]{Heston/images/srforgan_samples_2.png}
    \caption{Samples obtained using fixed $x_0$, $\mu$, $v_0$, $\theta$, $\kappa$, $\sigma_v$, $\rho$.}
    \label{srforgan_heston_sample2}
\end{figure}



To validate the distributional predictions of the generative model against the
theoretical Heston dynamics, an analytical consistency check is performed on the
conditional mean and standard deviation of the ten-step-ahead log-price
distribution.  The analysis compares three quantities: the theoretical moments
derived from the Heston characteristic function, the empirical moments of the
Monte Carlo ground-truth distribution, and the empirical moments of the
generated distribution.\\

Under the Heston model the log-price follows the Stochastic Differential Equation described in equation \ref{eq:logprice_sde_heston},
so that the exact conditional mean over a forecast horizon $\tau$ is
\begin{equation}
    \mathbb{E}\bigl[X_{T+\tau}\mid X_T,\,v_T\bigr]
    \;=\;
    X_T \;+\; \mu\tau \;-\; \tfrac{1}{2}\,\bar{v}\,\tau,
    \qquad
    \bar{v} = \frac{\mathbb{E}\!\left[\int_0^\tau v_{T+s}\,ds\,\Big|\,v_T\right]}{\tau},
    \label{eq:heston_mean}
\end{equation}
where the integrated variance is given in closed form by the CIR expectation
\begin{equation}
    \mathbb{E}\!\left[\int_0^\tau v_{T+s}\,ds\,\Big|\,v_T\right]
    \;=\;
    \theta\tau + (v_T - \theta)\,\frac{1 - e^{-\kappa\tau}}{\kappa}.
    \label{eq:integral_v}
\end{equation}
Since $\mu = 0$ in our parametrisation, the only deviation from $X_T$ is the
It\^{o} correction $-\tfrac{1}{2}\mathbb{E}[\int v_s\,ds]$, which is always
negative.  For the ten-step horizon $\tau = 10/252 \approx 0.040$\,yr with
$\kappa\tau = 0.079 \ll 1$, equation~\eqref{eq:integral_v} reduces to
$\approx v_T\tau$ with negligible error, so the mean shift from $X_T$ is
$-\tfrac{1}{2}v_T\tau$.

The theoretical standard deviation is dominated by the same integrated variance:
\begin{equation}
    \sigma_{\mathrm{th}} \;=\; \sqrt{\,\mathbb{E}\!\left[\int_0^\tau v_{T+s}\,ds\right]},
    \label{eq:std_th}
\end{equation}
which equals $\sqrt{v_T\tau}$ to leading order.  The vol-of-vol correction
term, proportional to $\sigma_v^2\tau^2$, contributes an additional
$\approx 12.5\%$ to the standard deviation across all four trajectories;
this small and uniform inflation reflects the brevity of the horizon relative
to the mean-reversion timescale.\\


Table~\ref{tab:mean_consistency} reports the theoretical mean from
equation~\eqref{eq:heston_mean}, together with the empirical means of the
true and generated Monte Carlo distributions.

\begin{table}[H]
\centering
\caption{Conditional mean consistency.  $\hat{\mu}_{\mathrm{true}}$ and
$\hat{\mu}_{\mathrm{gen}}$ are the empirical means of the Monte Carlo ground
truth and of the generated distribution, respectively.
$|\Delta_{\mathrm{true}}|=|\hat{\mu}_{\mathrm{true}}-\mu_{\mathrm{th}}|$,
$|\Delta_{\mathrm{gen}}|=|\hat{\mu}_{\mathrm{gen}}-\mu_{\mathrm{th}}|$.}
\label{tab:mean_consistency}
\begin{tabular}{cccccccc}
\toprule
Traj. & $X_T$ & $v_T$ & $\mu_{\mathrm{th}}$
      & $\hat{\mu}_{\mathrm{true}}$ & $|\Delta_{\mathrm{true}}|$
      & $\hat{\mu}_{\mathrm{gen}}$\phantom{$|$}
      & $|\Delta_{\mathrm{gen}}|$ \\
\midrule
0 & $-1.219$ & $0.879$ & $-1.236$ & $-1.237$ & $0.001$ & $-1.243$ & $0.006$\\
1 & $-0.975$ & $1.094$ & $-0.996$ & $-0.983$ & $0.014$ & $-1.086$ & $0.09$\\
2 & $\phantom{-}0.721$ & $0.367$  & $\phantom{-}0.714$ & $\phantom{-}0.712$ & $0.001$ & $\phantom{-}0.675$ & $0.039$ \\
3 & $-0.159$ & $0.498$ & $-0.169$ & $-0.174$ & $0.006$ & $-0.147$ & $0.022$\\
\bottomrule
\end{tabular}
\end{table}

The true distribution means match the theoretical predictions to within
$0.014$ in all cases, confirming that the Monte Carlo simulation correctly
implements the Heston dynamics.\\
A noteworthy feature of equation~\eqref{eq:heston_mean} is that, with $\mu=0$,
the distribution centre lies slightly \emph{left} of the current log-price
$X_T$ by the It\^{o} correction $-\tfrac{1}{2}v_T\tau \in [-0.021,\,-0.007]$.\\


Table \ref{tab:std_consistency} compares the theoretical standard deviation
from equation~\eqref{eq:std_th} against the empirical counterparts.

\begin{table}[H]
\centering
\caption{Conditional standard deviation consistency.
$\sigma_{\mathrm{th}} = \sqrt{\mathbb{E}[\int_0^\tau v_s\,ds]}$;
$\hat{\sigma}_{\mathrm{true}}$ and $\hat{\sigma}_{\mathrm{gen}}$ are the
empirical standard deviations of the true and generated distributions.}
\label{tab:std_consistency}
\begin{tabular}{cccccccc}
\toprule
Traj. & $v_T$ & $\sqrt{v_T\tau}$ & $\sigma_{\mathrm{th}}$
      & $\hat{\sigma}_{\mathrm{true}}$ & $|\Delta_{\mathrm{true}}|$
      & $\hat{\sigma}_{\mathrm{gen}}$ & $|\Delta_{\mathrm{gen}}|$ \\
\midrule
0 & $0.879$ & $0.1868$ & $0.1852$ & $0.1830$ & $0.0022$ & $0.2068$ & $0.0216$ \\
1 & $1.094$ & $0.2084$ & $0.2062$ & $0.2129$ & $0.0067$ & $0.1989$ & $0.0073$ \\
2 & $0.367$ & $0.1207$ & $0.1215$ & $0.1227$ & $0.0012$ & $0.1160$ & $0.0055$ \\
3 & $0.498$ & $0.1406$ & $0.1406$ & $0.1443$ & $0.0037$ & $0.1377$ & $0.0029$ \\
\bottomrule
\end{tabular}
\end{table}

The true empirical standard deviations deviate from the theoretical values
by at most $0.007$, well within Monte Carlo sampling noise at $10^3$
simulations.  Crucially, the \emph{ordering} imposed by $v_T$ is preserved:
trajectory~2 ($v_T = 0.367$, lowest variance) yields the narrowest
distribution and trajectory~1 ($v_T = 1.094$, highest) the widest, in both
the theoretical and empirical cases.  The generated standard deviations are also
in reasonable agreement with theory, with trajectory~0 exhibiting the largest
absolute deviation ($0.022$).\\


The analytical assessment confirms full consistency between the Heston
parameter set, the terminal states $(X_T,\,v_T)$, and the Monte Carlo ground-truth
distributions: theoretical means match empirical means to within $0.014$ and
theoretical standard deviations to within $0.007$.  The generated distributions
replicate the correct distributional shape and ordering as suggested by the 
distance metrics and Kolmogorov-Smirnov test as well.\\


The same kind of experiment is run on Heston dataset obtained by allowing 
the different parameters to change.\\
The first set of experiments allows just the volatility of the volatility to change. In this case the parameters of the trajectories are sampled as follows:
\begin{table}[H]
\centering
\begin{tabular}{lcc}
\hline
Parameter & Value training & Value testing \\
\hline
$x_0$      & 0.0 & 0.0\\
$\mu$      & 0.0 & 0.0\\
$v_0$      & 0.5 & 0.5\\
$\kappa$   & 2.0 & 2.0\\
$\theta$   & 0.5 & 0.5\\
$\sigma_v$ & (0.2, 1.3) & (0.3, 1.2) \\
$\rho$     & -0.7 & -0.7\\
\texttt{Condition\_trajectory\_length} & 252 & 252\\
\texttt{Forecast\_horizon}     & 10 & 10\\
\hline
\end{tabular}
\caption{Heston parameters vol-of-vol}
\end{table}

The performance metrics obtained over \text{10000} samples are slightly decreased
with respect to the fixed parameters experiment:

\begin{table}[h]
\centering
\begin{tabular}{lc}
\hline
Metric & Value \\
\hline
KS test acceptance rate & 83.00\% \\
Jensen--Shannon distance & 0.149905 \\
Hellinger distance & 0.128144 \\
Total variation distance & 0.128569 \\
Earth mover's distance & 0.001136 \\
\hline
\end{tabular}
\caption{Distribution comparison metrics}
\end{table}

\begin{figure}[H]
    \centering
    \includegraphics[width=1.0\textwidth]{Heston/images/srforgan_samples_volvol.png}
    \caption{Samples obtained using fixed $x_0$, $\mu$, $v_0$, $\theta$, $\kappa$, $\rho$ and variable $\sigma_v$}
    \label{srforgan_heston_sample_volvol}
\end{figure}

The last experiment allows all the different parameters to change for each trajectory,
while keeping the Feller condition always satisfied, 
in particular:

\begin{table}[H]
\centering
\begin{tabular}{lcc}
\hline
Parameter & Value training & Value testing \\
\hline
$x_0$      & 0.0  & 0.0 \\
$\mu$      & 0.0 & 0.0 \\
$v_0$      & 0.5 & 0.5 \\
$\kappa$   & (1.5, 5.0) & (2.0, 4.0) \\
$\theta$   & (0.57, 1.5) & (0.6, 1.4) \\
$\sigma_v$ & (0.2, 1.3) & (0.3, 1.2) \\
$\rho$     & -0.7 & -0.7 \\
\texttt{Condition\_trajectory\_length} & 252 & 252 \\
\texttt{Forecast\_horizon}     & 10 & 10\\
\hline
\end{tabular}
\caption{Heston with all variable parameters}
\end{table}


The performance metrics does not differ significantly from the previous experiment where just the vol-of-vol 
is allowed to change for each trajectory. This suggests that the model is able to 
model also trajectories guided by different underlying parameters 
\begin{table}[h]
\centering
\begin{tabular}{lc}
\hline
Metric & Value \\
\hline
KS test acceptance rate & 82.10\% \\
Jensen--Shannon distance & 0.150024 \\
Hellinger distance & 0.129049 \\
Total variation distance & 0.129361 \\
Earth mover's distance & 0.001198 \\
\hline
\end{tabular}
\caption{Distribution comparison metrics}
\end{table}



\begin{figure}[H]
    \centering
    \includegraphics[width=1.0\textwidth]{Heston/images/srforgan_samples_all_params.png}
    \caption{Samples obtained using fixed $x_0$, $\mu$, $v_0$, but different $\theta$, $\kappa$, $\rho$ and $\sigma_v$}
    \label{srforgan_heston_sample_all_var}
\end{figure}



\section{Pricing-Based Validation of the Forecast Distribution}
\label{sec:pricing_validation}


The distributional metrics and experiments ran before quantify how close 
the generated distribution $P^\phi(\cdot\mid\mathbf c)$
is to the true distribution$P^\star(\cdot\mid\mathbf c)$ in a generic, application-agnostic environment.\\
The following experiments in this section evaluate the SRForGAN generated distribution in a pricing setting by using European options.
If $G_\phi$ has genuinely learned the conditional law of
$S_{t+l}$, discounted expectations of option payoffs computed under
$P^\phi$ should agree with the corresponding Heston-model prices, across
a grid of strikes spanning the smile.\\

A single price comparison, may conflates two distinct sources of
discrepancy, that is, the sampling noise inherent to any Monte Carlo estimator of
an expectation, and genuine mismatch between $P^\phi$ and $P^\star$. To
disentangle them, three pricing engines are used side by side.
\begin{itemize}
    \item An \emph{oracle} price is obtained semi-analytically from the Heston
characteristic function using the Fourier inversion method already described in \ref{alg:frFTT} 
which is the closest available proxy to the true,
noise-free price, but one that requires knowledge of the latent
instantaneous variance $v_t$, which $G_\phi$ never observes.
    \item A \emph{Heston Monte Carlo} price is obtained with the same
discounted-payoff sample-average estimator as the GAN price, but drawing
from the true Heston dynamics started from the
same $(S_t,v_t)$ like in \ref{alg:heston_mc_pdf}.
    \item A \emph{GAN} price is obtained with the identical
sample-average estimator, drawing instead from $G_\phi(\cdot\mid\mathbf
c)$.
\end{itemize}
   
Because the Heston-MC and GAN routes share the same estimator
family, their standard errors are directly comparable and
$\text{price}^{\text{GAN}}-\text{price}^{\text{MC}}$ isolates the
generator's distributional error net of Monte Carlo noise; because the
oracle and Heston-MC routes both target the true Heston distribution,
$\text{price}^{\text{MC}}-\text{price}^{\text{oracle}}$ measures pure
numerical noise and acts as a control, in fact, a large gap here would flag a
problem with the benchmark itself rather than with $G_\phi$.

\subsection{Experimental design}

\paragraph{Test set.} $J=300$ fresh Heston trajectories (Milstein scheme)
are simulated with parameters drawn from the training ranges, that allow the Feller condition to be always satisfied, of the previous experiment,
$\kappa\in[1.5,5.0]$, $\theta\in[0.57,1.5]$, $\sigma_v\in[0.2,1.3]$, and
$\rho=-0.7$, $v_0=0.5$ fixed. Each
trajectory carries a conditioning window of $k=252$ observed log-price
steps ($\approx1$ trading year) and a forecast horizon of $l=10$ steps
($\approx2$ weeks), matching the training configuration. The reference
time $t$ is the last index of the conditioning window; $S_t=e^{X_t}$ and
$v_t$ (both known from simulation, but $v_t$ withheld from $G_\phi$)
anchor all three pricing engines to the same starting state.

\paragraph{Strike grid.} For every trajectory, strikes are built on a
common moneyness grid $\mathcal M=\{0.80,0.90,0.95,1.00,1.05,1.10,1.20\}$,
$K=m\cdot S_t$, so each trajectory contributes the same relative coverage
of the smile regardless of its own realised spot level. Following
standard smile-construction practice, options are priced as puts for
$K<S_t$ and calls for $K\ge S_t$, keeping every quoted contract
out-of-the-money and away from the numerically unstable deep-ITM region
of Black--Scholes implied-volatility inversion.

\begin{table}[H]
\centering
\begin{tabular}{cc}
\hline
$K/S_t$ & \textbf{Contract} \\
\hline
0.80 & Put \\
0.90 & Put \\
0.95 & Put \\
1.00 & Call \\
1.05 & Call \\
1.10 & Call \\
1.20 & Call \\
\hline
\end{tabular}
\end{table}



\paragraph{Discounting.} The risk-free rate is set to $r=0$, matching the
risk-neutral drift $\mu=0$ used throughout the simulated data, so $S_t$
is a martingale under both the true and the forecast measure and the
three engines remain directly comparable.

\paragraph{Oracle price.} Writing $\tau=l\Delta t$, the Heston
characteristic function $\varphi(u)=\mathbb E[e^{iuX_{t+\tau}}\mid
X_t,v_t]$ is inverted with a fractional FFT into a discretised density
$f(x)$ of the terminal log-price.
\[
C^{\text{oracle}}(K)=e^{-r\tau}\int (e^{x}-K)^+ f(x)\,dx .
\]

Similarly for the put options with $K<S_t$:

\[
P^{\text{oracle}}(K)=e^{-r\tau}\int (K-e^{x})^+ f(x)\,dx
\]

During the rest of the discussion the notation $\hat C^{\text{°}}(K)$ will be used to refer to both call and put options, with the understanding that the payoff function is chosen according to the moneyness of the strike $K$.


\paragraph{Heston Monte Carlo price.} $10{,}000$ terminal log-price draws
per trajectory are simulated with the exact/QE CIR scheme already used
in this thesis, started from $(X_t,v_t)$:
\[
\hat C^{\text{MC}}(K)=\frac{e^{-r\tau}}{N}\sum_{i=1}^N
\left(e^{X_{t+\tau}^{(i)}}-K\right)^{+},\qquad
\widehat{\text{se}}^{\text{MC}}(K)=\frac{e^{-r\tau}}{\sqrt N}\,
\widehat{\mathrm{sd}}\!\left[\left(e^{X_{t+\tau}}-K\right)^{+}\right].
\]


\paragraph{GAN price.} The identical estimator is applied to $N=10{,}000$
draws $\hat X_{t+\tau}^{(i)}\sim G_\phi(\cdot\mid\mathbf c)$ from the
trained generator, conditioned only on the observed price window.
Similarly to the Heston-MC case, the standard error of the GAN price is estimated as:
\[
\hat C^{\text{GAN}}(K)=\frac{e^{-r\tau}}{N}\sum_{i=1}^N
\left(e^{X_{t+\tau}^{(i)}}-K\right)^{+},\qquad
\widehat{\text{se}}^{\text{GAN}}(K)=\frac{e^{-r\tau}}{\sqrt N}\,
\widehat{\mathrm{sd}}\!\left[\left(e^{X_{t+\tau}}-K\right)^{+}\right].
\]

\paragraph{Implied volatility.} Every price (oracle, MC, GAN) is
converted into a Black--Scholes implied volatility $\sigma_{\text{impl}}$
by Brent root-finding. By using this monotone re-parametrisation the prices errors are placed
 on a scale comparable across strikes and that directly exposes smile-shape mismatches, which is the
feature a stochastic-volatility model is specifically built to capture.

\paragraph{Error and diagnostic metrics.} For each (trajectory, strike):
price and IV errors of GAN vs.\ oracle, GAN vs.\ MC, and MC vs.\ oracle;
and a $z$-score
testing whether the GAN/MC price gap exceeds sampling noise. Separately,
for each trajectory the first four empirical moments of the raw MC and
GAN log-price samples (mean, standard deviation, skewness, excess
kurtosis) are compared directly, to diagnose which feature of the
distribution drives any pricing mismatch.\\

Let $i=1,\dots,n$ ($n=2{,}100$) index the
(trajectory, strike) pairs underlying Table~\ref{tab:overall}. For pair
$i$, let $\hat C_i^{\text{oracle}}$, $\hat C_i^{\text{MC}}$,
$\hat C_i^{\text{GAN}}$ denote the oracle, Heston-MC, and GAN prices, and
$\hat\sigma_i^{\text{oracle}}$, $\hat\sigma_i^{\text{GAN}}$ the
corresponding Black--Scholes implied volatilities. Define the residuals
\[
e_i = \hat C_i^{\text{GAN}}-\hat C_i^{\text{oracle}}, \qquad
e_i^{\text{MC}} = \hat C_i^{\text{GAN}}-\hat C_i^{\text{MC}}, \qquad
\varepsilon_i = \hat\sigma_i^{\text{GAN}}-\hat\sigma_i^{\text{oracle}},
\]
and the standardised price gap
\[
z_i = \frac{\hat C_i^{\text{GAN}}-\hat C_i^{\text{MC}}}
{\sqrt{\big(\widehat{\text{se}}_i^{\text{GAN}}\big)^2+\big(\widehat{\text{se}}_i^{\text{MC}}\big)^2}},
\]
where $\widehat{\text{se}}_i^{\text{GAN}}$ and $\widehat{\text{se}}_i^{\text{MC}}$
are the Monte Carlo standard errors of the GAN and Heston-MC price
estimators defined in the "Heston Monte Carlo price"/"GAN price"
paragraphs above, and $\mathbf 1\{\cdot\}$ is the indicator function.

\subsection{Results}

Table~\ref{tab:overall} reports the aggregate error statistics over all
$J\times|\mathcal M|=2{,}100$ (trajectory, strike) pairs.
Table~\ref{tab:moneyness} breaks these down by moneyness, and
Table~\ref{tab:moments} compares the empirical moments of the GAN and
Heston-MC terminal distributions.\\



\begin{table}[H]
\centering
\small
\caption{Aggregate pricing-error statistics ($J=300$ trajectories,
$7$ strikes, $n=2{,}100$).}
\label{tab:overall}
\begin{tabular}{lcc}
\toprule
Metric & Formula & Value \\
\midrule
Bias, price vs.\ oracle          & $\dfrac1n\sum_i e_i$                                   & 0.0005 \\[1.2ex]
\hline
RMSE, price vs.\ oracle          & $\sqrt{\dfrac1n\sum_i e_i^2}$                         & 0.0312 \\[1.2ex]
\hline
WAPE, price vs.\ oracle           & $\frac{\sum \vert{}e_i\vert{}}{\sum C_i^{\text{oracle}}}$  & 0.3718 \\[1.2ex]
\hline
RMSE, price vs.\ Heston MC       & $\sqrt{\dfrac1n\sum_i \big(e_i^{\text{MC}}\big)^2}$    & 0.0311 \\[1.2ex]
\hline
WAPE, price vs.\ Heston MC        & $\frac{\sum \vert{}e_i^{MC}\vert{}}{\sum C_i^{\text{oracle}}}$             & 0.3731 \\[1.2ex]
\hline
Bias, implied vol.\ vs.\ oracle  & $\dfrac1n\sum_i \varepsilon_i$                         & 0.0045 \\[1.2ex]
\hline
RMSE, implied vol.\ vs.\ oracle  & $\sqrt{\dfrac1n\sum_i \varepsilon_i^2}$                & 0.3320 \\[1.2ex]
\hline
WAPE, implied vol.\ vs.\ oracle   & $\frac{\sum \vert{}\varepsilon_i\vert{}}{\sum \hat\sigma_i^{\text{oracle}}}$                      & 0.2601 \\[1.2ex]
\hline
Share of $|z_{\text{GAN vs.\ MC}}|>1.96$ & $\dfrac1n\sum_i \mathbf 1\{|z_i|>1.96\}$       & 91.3\% \\
\bottomrule
\end{tabular}
\end{table}


\begin{table}[H]
\centering
\caption{Pricing-error statistics by moneyness $K/S_t$. \\(*** = $p<0.01$, ** = $p<0.05$, * = $p<0.1$).}
\label{tab:moneyness}
\begin{tabular}{lrrrr}
\toprule
$K/S_t$ & Bias, price & RMSE, IV & WAPE, IV & Bias, IV \\
\midrule
0.80 & $-0.0023$\text{***} & 0.2256 & 0.4696 & $-0.0574$\text{***}\\
0.90 & $-0.0041$\text{***} & 0.2861 & 0.3494 & $-0.0534$\text{**} \\
0.95 & $-0.0048$\text{**} & 0.3300 & 0.2961 & $-0.0517$\text{**} \\
1.00 & $\phantom{-}0.0049$\text{*} & 0.2688 & 0.3557 & $\phantom{-}0.0662$\text{*} \\
1.05 & $\phantom{-}0.0042$\text{**} & 0.3133 & 0.2951 & $\phantom{-}0.0541$\text{*} \\
1.10 & $\phantom{-}0.0035$\text{**} & 0.3679 & 0.2522 & $\phantom{-}0.0432$\text{**} \\
1.20 & $\phantom{-}0.0021$\text{***} & 0.4851 & 0.1985 & $\phantom{-}0.0303$\text{***} \\
\bottomrule
\end{tabular}
\end{table}

\begin{table}[H]
\centering
\caption{GAN vs.\ Heston-MC terminal log-price moments (bias
$=\text{GAN}-\text{MC}$, averaged over $J=300$ trajectories; correlation
across trajectories).}
\label{tab:moments}
\begin{tabular}{lrrr}
\toprule
Moment & Bias & WAPE & Correlation \\
\midrule
Mean               & 0.0095  & 0.0516 & 0.998 \\
Std.\ deviation     & $-0.0103$ & 0.101 & 0.814 \\
Skewness            & 0.0725  & 0.4964 & $-0.099$ \\
Excess kurtosis     & 0.1078  & 1.5267 & 0.374 \\
\bottomrule
\end{tabular}
\end{table}

The oracle and Heston-MC benchmarks agree closely with each other, both
in aggregate (RMSE $0.0312$ vs.\ $0.0311$) and at every moneyness level;
this confirms the two ground-truth routes are mutually consistent, so
essentially all the reported GAN discrepancy reflects genuine
distributional error rather than benchmark disagreement. The aggregate
price bias is small ($0.0005$), but Table~\ref{tab:moneyness} shows this
is the average of a sign-flipping pattern: prices are systematically
\emph{below} the oracle for $K<S_t$ (puts) and \emph{above} it for
$K\ge S_t$ (calls), with the effect largest near the money. For $91.3\%$
of (trajectory, strike) pairs the GAN/MC price gap exceeds what sampling
noise alone would explain at the $95\%$ level ($|z|>1.96$), so the
mismatch is not an artefact of Monte Carlo variance. 
Moreover, for each moneyness level the bias in the prices and in the implied volatilities is statistically significant at least at the $10\%$ confidence level.\\

Table~\ref{tab:moments} identifies the mechanism. The GAN reproduces the
trajectory-conditional \emph{mean} of the terminal log-price almost
perfectly (correlation $0.998$) and the \emph{standard deviation}
reasonably well (correlation $0.814$), but shows essentially no
association with the true \emph{skewness} ($-0.099$) or \emph{kurtosis}
($0.374$) across trajectories. A small positive mean bias ($+0.0095$,
under $10\%$ of one standard deviation at this horizon) shifts the whole
distribution to the right, mechanically depressing put prices and
inflating call prices, exactly the sign pattern observed. The
Heston leverage effect ($\rho=-0.7$) induces genuine, trajectory-specific
negative skew in $P^\star$; since the GAN does not track skewness
across regimes, its implied smile comes out flatter and less negatively
skewed than the truth, which is precisely why implied-volatility errors
are largest at the money (where price is most sensitive to dispersion
and to skew-driven asymmetry) and shrink towards both wings, rather than
following the pattern of noisier deep-OTM/ITM inversion that one might
otherwise expect. A plausible contributing cause is that Energy-Score
training here used only $n_{\text{samples}}=10$ generator draws per
gradient step: mean and variance are estimated far more precisely than
skewness or kurtosis at that sample size, so the training signal is
inherently weaker for exactly the moments found deficient here. 
In addition, no hyperparameter tuning was performed to improve the GAN's ability to capture higher-order moments.



\chapter{Conclusion}

This thesis proposes to test and evaluate whether generative adversarial networks,
and their adversarial-free variant, can be used to learn the conditional distribution of a financial time series.
The proposed scoring rule based final model, named SRForGAN, has been tested to check its capabilities to learn the
forecast distribution of a financial time series accurately enough to be
useful, moving progressively from Geometric Brownian Motion to the
stochastic-volatility Heston model, and ultimately evaluating the learned
distributions in a concrete option pricing task.\\

The results consistently favour the proposed approach over its
GAN-based alternatives. Under Black--Scholes dynamics, the vanilla ForGAN
trained with the BCE loss struggled visibly (25\% KS acceptance,
Jensen--Shannon distance $\approx0.49$), the Wasserstein variant improved
matters but remained far from satisfactory (34\% KS acceptance,
JS $\approx0.28$), while SRForGAN, trained by minimising the Energy
Score, with no critic and no adversarial game, matched the analytical
Gaussian benchmark almost exactly (99\% KS acceptance, JS $\approx0.14$).
This gap is the central practical advantage demonstrated throughout the
thesis: replacing the adversarial min-max objective with a strictly proper
scoring rule removes the well-known instability and mode-collapse
tendencies of GAN training, at no cost in distributional accuracy and, in
this setting, at a considerable gain. The same architecture and training
recipe then transferred, with only a modest and expected drop in
performance, to the substantially harder Heston setting: 87.5\% KS
acceptance with fixed parameters, and 82--83\% once $\kappa$, $\theta$,
$\sigma_v$ and $\rho$ were all allowed to vary across trajectories, showing
that a single trained generator can reasonably generalise across a family
of volatility regimes rather than requiring re-estimation for each one,
which is precisely the property one would want from a model intended for
practical deployment. The machine-epsilon-level validation of the frFFT
inversion pipeline further ensures that all of these comparisons rest on a
numerically trustworthy Heston benchmark rather than on an artefact of the
evaluation code itself.

The pricing-based validation introduced in the final chapter is arguably
the thesis's most informative contribution, precisely because it complicates
this otherwise positive picture. Distributional distances such as KS,
Jensen--Shannon, Hellinger and Total Variation are dominated by the bulk of
the distribution and, as such, rewarded SRForGAN generously. Pricing a
volatility smile, by contrast, is disproportionately sensitive to exactly
the features that live in the tails and in the asymmetry of the
distribution. The moment diagnostics make the mechanism explicit: the
generator reproduces the trajectory-conditional mean of the terminal
log-price almost perfectly (correlation $0.998$) and the standard deviation
well (correlation $0.814$), but shows essentially no ability to track
skewness ($-0.189$) or excess kurtosis ($0.264$) across the different Heston
regimes. Since the leverage effect ($\rho=-0.7$) is precisely what endows
the true distribution with a regime-dependent negative skew, its absence in
the generated distribution produces a smile that is systematically flatter
than the truth, with a statistically significant bias (at worst $p<0.001$)
that underprices out-of-the-money puts and overprices out-of-the-money
calls. In other words, a model that looks close to indistinguishable from
the ground truth under generic distributional metrics can still be
economically unsuitable for a skew-sensitive financial application, a
distinction that only becomes visible once the forecast is stress-tested
against the task it is meant to serve, and one that future generative
forecasting work in finance should not overlook.

This shortcoming also delineates the clearest avenues for future work.
Training used only ten Monte Carlo draws per gradient step to estimate the
Energy Score, a sample size sufficient to estimate a mean or a variance
precisely but a considerably weaker basis for estimating skewness or
kurtosis; whether increasing this sample size, or adopting a scoring rule
more explicitly sensitive to higher-order moments (for instance a variogram
score or a weighted combination of scoring rules), closes this gap is a
natural and immediately testable next step, since no hyperparameter tuning
was directed at this issue in the present work. Likewise, the generator was
deliberately restricted to conditioning on the observed price history
alone, mirroring realistic deployment where the instantaneous variance is
unobservable; enriching the conditioning set with realised-volatility
proxies, extending the analysis to longer horizons and more extreme
parameter regimes, and testing the same pricing-based methodology on other
volatility models are all promising directions building directly on the
framework developed here.